% GENERATED FILE. Edit canonical lessons, then run npm run generate:books.
% canonical-source-hash: fnv1a64:f980a601b8d91b1c
% canonical-lessons: JA-C379-suwaru, JA-C379-sharyou, JA-C379-tsuuro, JA-C379-shanai, JA-C379-unchin

\chapter{Finding a Seat}
\label{ch:ja-finding-a-seat}
\hlchaptermodality{379}

\begin{chapteropening}
\textbf{By the end of this chapter:} \emph{I can find a seat and talk about the carriage.}

Complete the last lesson of chapter 379: I can find a seat and talk about the carriage.
\end{chapteropening}

\section[\texorpdfstring{\ja{すわる} (suwaru)}{すわる (suwaru)}]{\ja{すわる} (suwaru) — to sit down}

\label{lesson:JA-C379-suwaru}



\begin{quote}
Before the new one: say the Japanese for to cram in, then the Japanese for to stand up.
\end{quote}

\subsection*{You'll want to know: \ja{すわる}}
\textbf{\ja{すわる}} — \emph{suwaru} — \textquotedblleft{}to sit down\textquotedblright{}.

\begin{quote}
\textbf{\ja{でんしゃで} \ja{ははの} \ja{となりに} \ja{すわりました。}} — \emph{densha de haha no tonari ni suwarimashita} — \textquotedblleft{}On the train I sat next to my mother.\textquotedblright{}
\end{quote}

Polite \textbf{\ja{すわります}}, \emph{suwarimasu}; a godan verb. Past (polite): \textbf{\ja{おじいさんが} \ja{ざせきに} \ja{すわりました。}}, \emph{ojiisan ga zaseki ni suwarimashita}.

\begin{cousinweb}
\textbf{\ja{すわる}} (\emph{suwaru}) is a native Japanese word (wago), inherited rather than borrowed from Chinese; its kanji was attached later as a kun'yomi (JMdict).
\end{cousinweb}

\begin{grammarlens}[title={What you've built}]
One more word for getting around.
\end{grammarlens}

\subsection*{Your turn}
\begin{itemize}
\raggedright
  \item \emph{Say it:} \emph{suwaru}
  \item \emph{Say it:} \emph{suwaru}, once more
  \item \emph{Say it:} \emph{densha de haha no tonari ni suwarimashita} — On the train I sat next to my mother.
  \item \emph{Recall:} say the Japanese for to cram in, then the Japanese for to stand up, then say \emph{suwaru} again
\end{itemize}

\begin{tcolorbox}[breakable,colback=teal!4,colframe=teal!35!black,title={Before you move on}]
What does \ja{すわる} mean? (\textquotedblleft{}To sit down\textquotedblright{}.) Say it once more.
\end{tcolorbox}
\section[\texorpdfstring{\ja{しゃりょう} (sharyō)}{しゃりょう (sharyō)}]{\ja{しゃりょう} (sharyō) — a railway carriage}

\label{lesson:JA-C379-sharyou}



\begin{quote}
Before the new one: say the Japanese for to stand up, then the Japanese for to sit down.
\end{quote}

\subsection*{You'll want to know: \ja{しゃりょう}}
\textbf{\ja{しゃりょう}} — \emph{sharyō} — \textquotedblleft{}a railway carriage\textquotedblright{}.

\begin{quote}
\textbf{\ja{まえの} \ja{しゃりょうは} \ja{すいていました。}} — \emph{mae no sharyō wa suite imashita} — \textquotedblleft{}The front carriage was empty.\textquotedblright{}
\end{quote}

\begin{cousinweb}
\textbf{\ja{しゃりょう}} (\emph{sharyō}) is Sino-Japanese (kango): two Chinese characters read with their on'yomi, readings borrowed from Chinese. They are \emph{sha}, \textquotedblleft{}car\textquotedblright{} (Middle Chinese \emph{tsyhae}) and \emph{ryō}, \textquotedblleft{}counter for carriages\textquotedblright{} (\emph{ljangH}) (KANJIDIC2; Middle Chinese in Baxter's transcription).
\end{cousinweb}

\begin{grammarlens}[title={What you've built}]
One more word for getting around.
\end{grammarlens}

\subsection*{Your turn}
\begin{itemize}
\raggedright
  \item \emph{Say it:} \emph{sharyō}
  \item \emph{Say it:} \emph{sharyō}, once more
  \item \emph{Say it:} \emph{mae no sharyō wa suite imashita} — The front carriage was empty.
  \item \emph{Recall:} say the Japanese for to stand up, then the Japanese for to sit down, then say \emph{sharyō} again
\end{itemize}

\begin{tcolorbox}[breakable,colback=teal!4,colframe=teal!35!black,title={Before you move on}]
What does \ja{しゃりょう} mean? (\textquotedblleft{}A railway carriage\textquotedblright{}.) Say it once more.
\end{tcolorbox}
\section[\texorpdfstring{\ja{つうろ} (tsūro)}{つうろ (tsūro)}]{\ja{つうろ} (tsūro) — an aisle; a passage}

\label{lesson:JA-C379-tsuuro}



\begin{quote}
Before the new one: say the Japanese for to sit down, then the Japanese for a railway carriage.
\end{quote}

\subsection*{You'll want to know: \ja{つうろ}}
\textbf{\ja{つうろ}} — \emph{tsūro} — \textquotedblleft{}an aisle; a passage\textquotedblright{}.

\begin{quote}
\textbf{\ja{つうろに} \ja{にもつを} \ja{おかないで} \ja{ください。}} — \emph{tsūro ni nimotsu o okanaide kudasai} — \textquotedblleft{}Please do not leave your bags in the aisle.\textquotedblright{}
\end{quote}

\begin{cousinweb}
\textbf{\ja{つうろ}} (\emph{tsūro}) is Sino-Japanese (kango): two Chinese characters read with their on'yomi, readings borrowed from Chinese. They are \emph{tsū}, \textquotedblleft{}pass through\textquotedblright{} (Middle Chinese \emph{thuwng}) and \emph{ro}, \textquotedblleft{}path\textquotedblright{} (\emph{luH}) (KANJIDIC2; Middle Chinese in Baxter's transcription).
\end{cousinweb}

\begin{grammarlens}[title={What you've built}]
One more word for getting around.
\end{grammarlens}

\subsection*{Your turn}
\begin{itemize}
\raggedright
  \item \emph{Say it:} \emph{tsūro}
  \item \emph{Say it:} \emph{tsūro}, once more
  \item \emph{Say it:} \emph{tsūro ni nimotsu o okanaide kudasai} — Please do not leave your bags in the aisle.
  \item \emph{Recall:} say the Japanese for to sit down, then the Japanese for a railway carriage, then say \emph{tsūro} again
\end{itemize}

\begin{tcolorbox}[breakable,colback=teal!4,colframe=teal!35!black,title={Before you move on}]
What does \ja{つうろ} mean? (\textquotedblleft{}An aisle; a passage\textquotedblright{}.) Say it once more.
\end{tcolorbox}
\section[\texorpdfstring{\ja{しゃない} (shanai)}{しゃない (shanai)}]{\ja{しゃない} (shanai) — inside the train}

\label{lesson:JA-C379-shanai}



\begin{quote}
Before the new one: say the Japanese for a railway carriage, then the Japanese for an aisle.
\end{quote}

\subsection*{You'll want to know: \ja{しゃない}}
\textbf{\ja{しゃない}} — \emph{shanai} — \textquotedblleft{}inside the train\textquotedblright{}.

\begin{quote}
\textbf{\ja{しゃないで} \ja{でんわを} \ja{しないで} \ja{ください。}} — \emph{shanai de denwa o shinaide kudasai} — \textquotedblleft{}Please do not make phone calls on the train.\textquotedblright{}
\end{quote}

\begin{cousinweb}
\textbf{\ja{しゃない}} (\emph{shanai}) is Sino-Japanese (kango): two Chinese characters read with their on'yomi, readings borrowed from Chinese. They are \emph{sha}, \textquotedblleft{}car\textquotedblright{} (Middle Chinese \emph{tsyhae}) and \emph{nai}, \textquotedblleft{}inside\textquotedblright{} (\emph{nwojH}) (KANJIDIC2; Middle Chinese in Baxter's transcription).
\end{cousinweb}

\begin{grammarlens}[title={What you've built}]
One more word for getting around.
\end{grammarlens}

\subsection*{Your turn}
\begin{itemize}
\raggedright
  \item \emph{Say it:} \emph{shanai}
  \item \emph{Say it:} \emph{shanai}, once more
  \item \emph{Say it:} \emph{shanai de denwa o shinaide kudasai} — Please do not make phone calls on the train.
  \item \emph{Recall:} say the Japanese for a railway carriage, then the Japanese for an aisle, then say \emph{shanai} again
\end{itemize}

\begin{tcolorbox}[breakable,colback=teal!4,colframe=teal!35!black,title={Before you move on}]
What does \ja{しゃない} mean? (\textquotedblleft{}Inside the train\textquotedblright{}.) Say it once more.
\end{tcolorbox}
\section[\texorpdfstring{\ja{うんちん} (unchin)}{うんちん (unchin)}]{\ja{うんちん} (unchin) — a fare}

\label{lesson:JA-C379-unchin}



\begin{quote}
Before the new one: say the Japanese for an aisle, then the Japanese for inside the train.
\end{quote}

\subsection*{You'll want to know: \ja{うんちん}}
\textbf{\ja{うんちん}} — \emph{unchin} — \textquotedblleft{}a fare\textquotedblright{}.

\begin{quote}
\textbf{\ja{みなとまでの} \ja{うんちんは} \ja{いくらですか。}} — \emph{minato made no unchin wa ikura desu ka} — \textquotedblleft{}How much is the fare to the harbour?\textquotedblright{}
\end{quote}

\begin{cousinweb}
\textbf{\ja{うんちん}} (\emph{unchin}) is Sino-Japanese (kango): two Chinese characters read with their on'yomi, readings borrowed from Chinese. They are \emph{un}, \textquotedblleft{}carry\textquotedblright{} (Middle Chinese \emph{hjunH}) and \emph{chin}, \textquotedblleft{}fare\textquotedblright{} (\emph{nrimH}) (KANJIDIC2; Middle Chinese in Baxter's transcription).
\end{cousinweb}

\begin{grammarlens}[title={What you've built}]
One more word for getting around.
\end{grammarlens}

\subsection*{Your turn}
\begin{itemize}
\raggedright
  \item \emph{Say it:} \emph{unchin}
  \item \emph{Say it:} \emph{unchin}, once more
  \item \emph{Say it:} \emph{minato made no unchin wa ikura desu ka} — How much is the fare to the harbour?
  \item \emph{Recall:} say the Japanese for an aisle, then the Japanese for inside the train, then say \emph{unchin} again
\end{itemize}

\begin{tcolorbox}[breakable,colback=teal!4,colframe=teal!35!black,title={Before you move on}]
What does \ja{うんちん} mean? (\textquotedblleft{}A fare\textquotedblright{}.) Say it once more.
\end{tcolorbox}
